\paragraph{Faithfulness in the manuscript's Noetherian coherent refinement.}

\begin{proposition}\label{prop:DE-E2-coherent-faithfulness}
\status{Proved; root and reduction-descent full independent proof/source audits passed.}
Let $S$ be locally Noetherian and let $X\to S$ be flat and locally of finite presentation.  Let $\mathcal E,\mathcal F$ be quasi-coherent crystals on the ordinary infinitesimal site of $X/S$, with $\mathcal F_X$ coherent.  Then evaluation on $X$ is injective on morphisms:
\[
 \Hom_{(X/S)_{\rm inf}}(\mathcal E,\mathcal F)
       \longrightarrow \Hom_{\OO_X}(\mathcal E_X,\mathcal F_X).
\]
In particular the restriction functor from coherent crystals to ordinary coherent stratified modules is faithful.  The assertion holds for the small and the big infinitesimal site, including when arbitrary, possibly non-Noetherian, test schemes are allowed.  No finiteness assumption on $\mathcal E$ is required.
\end{proposition}
\begin{proof}
Let $u:\mathcal E\to\mathcal F$ have $u_X=0$.  We first work over an affine base $S=\Spec R$ and an affine open $V=\Spec A\subset X$ with $A$ a finitely presented flat $R$-algebra.  Consider an affine Noetherian test $U=\Spec C\hookrightarrow T=\Spec B$, whose structural map $U\to X$ factors through $V$.  Thus $B\twoheadrightarrow C$ has a nilpotent kernel $I$, and there is an $R$-algebra map $A\to C$.  In the small site, after restricting to affine opens, we may take $U=V$.

The crystal property identifies $\mathcal F_T/I\mathcal F_T$ with the pullback of $\mathcal F_X$ to $U$, which is a finite $C$-module.  Lifting a finite set of generators and iterating $I^N=0$ shows that $\mathcal F_T$ is a finite $B$-module; this elementary nilpotent form of Nakayama's argument does not require $\mathcal F_T$ to be finite in advance.

Set
\[
 D=A\otimes_R B,\qquad
 J=\ker\bigl(D\longrightarrow C,\ a\otimes b\longmapsto a\overline b\bigr).
\]
The map to $C$ is surjective.  The ring $D$ is Noetherian, since it is a finitely presented $B$-algebra, and it is flat over $B$, since $A$ is flat over $R$.  For every $n\geq1$, the closed immersion
\[
 U\hookrightarrow T_n=\Spec(D/J^n)
\]
is an infinitesimal test, with compatible maps to $T$ and to $V$.  The latter map and $u_X=0$ give $u_{T_n}=0$.  The former map and the crystal transition isomorphisms therefore give
\[
 (D/J^n)\otimes_B u_T=0\qquad(n\geq1).
\]
Put $N=D\otimes_B\mathcal F_T$, a finite $D$-module.  Every element in the image of $D\otimes_B u_T$ belongs to $\bigcap_nJ^nN$.

Fix any prime $\mathfrak q\subset B$.  Since $I$ is nilpotent, it determines a prime of $C=B/I$ and hence a prime $\mathfrak p\subset D$ containing $J$, with $\mathfrak p\cap B=\mathfrak q$.  Krull's intersection theorem in the Noetherian local ring $D_{\mathfrak p}$ gives
\[
 \bigcap_{n\geq1}(JD_{\mathfrak p})^nN_{\mathfrak p}=0.
\]
Thus $D_{\mathfrak p}\otimes_Bu_T=0$.  The map $B_{\mathfrak q}\to D_{\mathfrak p}$ is flat and local, hence faithfully flat.  Consequently $(u_T)_{\mathfrak q}=0$.  As $\mathfrak q$ was arbitrary, $u_T=0$.  The precise commutative-algebra inputs are Krull intersection and faithful flatness of a flat local map, \citet[Lemmas~10.51.4 and~10.39.17, Tags~00IP and~00HR]{StacksProject}; the slightly more general separation statement is Lemma~10.51.5, Tag~00IQ.

It remains to remove the Noetherian condition on the test.  Locally write $A=R[z_1,\ldots,z_r]/K$, lift the images of the $z_i$ from $C$ to $B$, and choose $N$ with $I^N=0$.  The resulting map $R[z_1,\ldots,z_r]\to B$ annihilates $K^N$, and hence factors through the Noetherian ring
\[
 B_0=R[z_1,\ldots,z_r]/K^N.
\]
The test $(U,T)$ maps to the finite-type test
$V\hookrightarrow\Spec B_0$.  In the small site $U$ is an open of $V$, as required; in the big site its given map to $V$ suffices.  We already proved $u_{\Spec B_0}=0$.  The crystal transition isomorphisms then give $u_T=0$.  These local conclusions glue over $X$ and $S$, proving the proposition.
\end{proof}

\begin{remark}[What this proves in the coherent refinement]\label{rem:DE-E2-coherent-faithfulness-scope}
\status{Proved consequence of Proposition~\ref{prop:DE-E2-coherent-faithfulness}.}
The nonfaithfulness in Theorem~\ref{thm:DE-E2-counterexample} uses a non-Noetherian base, and that in Corollary~\ref{cor:DE-E2-Noetherian-counterexample} uses a noncoherent target.  The proposition shows that both cannot be avoided simultaneously.  The proposition itself asserts faithfulness, even with an arbitrary quasi-coherent source.  Fullness and essential surjectivity for coherent source and target are established separately in Theorem~\ref{thm:DE-E2-coherent-Noetherian}, using completed ordinary descent and an explicit zeroth-cohomology comparison.  A derived comparison theorem alone would not establish those ordinary conclusions.
\end{remark}
